\documentclass[fontsize=11pt, parskip=half]{scrartcl}
\usepackage[paperwidth=6.5in, paperheight=9.5in,
            margin=1cm]{geometry}
\pagestyle{plain}

\usepackage{amsmath, amssymb, amsthm, amsfonts}

\usepackage[dvipsnames]{xcolor}

\usepackage[draft]{commenting}
\declareauthor{leon}{LL}{blue}
\authorcommand{leon}{comment}

\usepackage[colorlinks=true, linkcolor=blue, urlcolor=blue, citecolor=blue]{hyperref}
\usepackage{url}
\usepackage[capitalise,noabbrev]{cleveref}

% 3. Define your theorem environments AFTER cleveref is loaded
\newtheorem{theorem}{Theorem}[section]
\newtheorem{lemma}[theorem]{Lemma}
\newtheorem{fact}[theorem]{Fact}
\theoremstyle{definition}
\newtheorem{definition}[theorem]{Definition}
\newtheorem{remark}[theorem]{Remark}

% 4. Set up cross-reference names
\crefname{section}{Problem}{Problems}
\Crefname{section}{Problem}{Problems}

\crefname{theorem}{Theorem}{Theorems}
\Crefname{theorem}{Theorem}{Theorems}

\crefname{definition}{Definition}{Definitions}
\Crefname{definition}{Definition}{Definitions}

\crefname{lemma}{Lemma}{Lemmas}
\Crefname{lemma}{Lemma}{Lemmas}

\crefname{fact}{Fact}{Facts}
\Crefname{fact}{Fact}{Facts}

% Solution toggle: set to true to show solutions, false to hide
\newif\ifsolutions
\solutionstrue  % comment out or set \solutionsfalse to hide

\usepackage{tcolorbox}
\tcbuselibrary{breakable}

\ifsolutions
  \newenvironment{solution}{%
    \begin{tcolorbox}[colback=black!5!white, colframe=black!50!white,
      title={\textbf{Solution}}, fonttitle=\normalfont, breakable]
  }{\end{tcolorbox}}
\else
  \newenvironment{solution}{\expandafter\comment}{\expandafter\endcomment}
  \usepackage{verbatim} % for \comment environment
\fi

\tcolorboxenvironment{definition}{colback=blue!3!white, colframe=blue!40!white,
  breakable, top=4pt, bottom=4pt}
\tcolorboxenvironment{theorem}{colback=red!3!white, colframe=red!40!white,
  breakable, top=4pt, bottom=4pt}
\tcolorboxenvironment{lemma}{colback=purple!3!white, colframe=purple!40!white,
  breakable, top=4pt, bottom=4pt}
\tcolorboxenvironment{fact}{colback=orange!3!white, colframe=orange!40!white,
  breakable, top=4pt, bottom=4pt}

\newtcolorbox{defbox}[1][]{colback=blue!3!white, colframe=blue!40!white,
  title={\textbf{#1}}, fonttitle=\normalfont, breakable, top=4pt, bottom=4pt}

\newtcolorbox{examplebox}[1][]{colback=green!3!white, colframe=green!40!black,
  title={\textbf{#1}}, fonttitle=\normalfont, breakable, top=4pt, bottom=4pt}
%%%%%%%%%%%%%%%%%%%%%%%%%

\newcommand{\mytitle}[1]{{\noindent\Large\textbf{#1}}}

\renewcommand{\thesection}{\arabic{section}}

\newcommand{\mysubsection}[2]{\par\medskip\noindent\refstepcounter{subsection}\textbf{\thesubsection.}~#2\par\smallskip}

\newcommand{\Msol}{\mathcal{M}_{sol}}

%===================================
\begin{document}

\noindent\textbf{Iliad Intensive} \hfill \textbf{London Initiative for Safe AI}\\
April 2026 \hfill Leon Lang\\

\mytitle{Solomonoff Induction Exercises \hfill \today}

\medskip
\noindent
\textbf{Difficulty ratings.} Each problem is tagged with a rating in square brackets following Knuth's exercise rating scheme: roughly, \textbf{[00]}~trivial, \textbf{[10]}~15--minute pencil-and-paper, \textbf{[20]}~1--2~hours, \textbf{[30]}~several hours to a day, \textbf{[40]}~a significant research result. Intermediate values are possible. See \cref{sec:appendix_difficulty} for the full scale.

%===================================
\section*{Goal and Roadmap}

\textbf{A recipe for prediction:} Solomonoff induction is an attempt to mathematically
formalize the \href{https://en.wikipedia.org/wiki/Problem_of_induction}\emph{problem of induction}
from philosophy: How to make predictions about the future based on past observations?

We formalize this as sequence prediction: There is some true environment (measure) $\mu$ that generates 
a sequence $x_1, x_2, \dots$ of (binary) symbols. Each next symbol $x_t \sim \mu(\cdot \mid x_{<t})$ is sampled
from $\mu$ conditioned on the past $x_{<t}$. A predictor $P$ takes the history $x_{<t}$ and gives a 
distribution $P(\cdot \mid x_{<t})$ over the next symbol $x_t$. 

We measure the quality of a predictor $P$ by the expected squared prediction error:
\begin{align}
S_\infty^\mu := \sum_{t=1}^{\infty} \sum_{x_{<t} \in \mathbb{B}^*} \mu(x_{<t}) \sum_{x_t \in \mathbb{B}} \big( P(x_t \mid x_{<t}) - \mu(x_t \mid x_{<t}) \big)^2.
\end{align}

Bayesian inference to the rescue: We consider a reference set $\mathcal{M}$ of possible environments,
allow our predictor $P$ to be a Bayesian mixture $\xi$ over the reference set,
$$
\xi(x_t \mid x_{<t}) := \sum_{\nu \in \mathcal{M}} w(\nu | x_{<t}) \nu(x_t \mid x_{<t}),
$$
where 
\begin{itemize}
\item $\nu(x_t \mid x_{<t})$ is how likely it is that $x_t$ would follow $x_{<t}$ if the sequence was
generated by $\nu$, and
\item $w(\nu | x_{<t})$ is how likely it is we believe $\nu$ to be the true environment,
given $x_{<t}$ was observed. 
\end{itemize}
We measure the quality of $\xi$ by the expected squared prediction error:
\begin{align}
S_\infty^\xi := \sum_{t=1}^{\infty} \sum_{x_{<t} \in \mathbb{B}^*} \xi(x_t \mid x_{<t}) \sum_{x_t \in \mathbb{B}} \big( \xi(x_t \mid x_{<t}) - \mu(x_t \mid x_{<t}) \big)^2.
\end{align}

We will prove results for a generic mixture $\xi$, and then later specialize to the universal mixture $\xi_U$:
%There are two parameters to consider: the reference class $\mathcal{M}$ and the prior $w(\cdot | x_{<t})$.

% \begin{itemize}
% \item We choose as $\mathcal{M}$ 
% the widest reasonable class: $\mathcal{M}_{sol}$, the set of all \emph{lower semicomputable semimeasures}, and as the prior 
% \item For the prior $w_\nu$ the Solomonoff prior $w_\nu = 2^{-K(\nu)}$, where $K(\nu)$ is the Kolmogorov complexity of $\nu$
% (length of the shortest program that computes $\nu$). 
% \end{itemize}
% We gloss over the details of lower semicomputability and Kolomogorov compel

% \textbf{Goal.} A Solomonoff predictor knows nothing about which world it lives in: rather than committing to a model, it hedges over \emph{every} semimeasure at once through the universal mixture $M$. This sheet proves that this single, universal predictor nevertheless learns to predict \emph{any} computable environment essentially as well as if it had known that environment all along. This is the central justification for taking $M$ as an idealized model of induction.

% \textbf{Main result (Solomonoff's prediction theorem).} Let $\mu$ be the true measure generating the data and let $M$ be the Solomonoff mixture. Define the \textbf{cumulative expected squared prediction error}
% \[
% S_\infty^\mu := \sum_{t=1}^{\infty} \sum_{x_{<t} \in \mathbb{B}^*} \mu(x_{<t}) \sum_{x_t \in \mathbb{B}} \big( M(x_t \mid x_{<t}) - \mu(x_t \mid x_{<t}) \big)^2.
% \]
% We prove that $S_\infty^\mu$ is \textbf{finite}: it is bounded by $-\ln w_\mu$ (\cref{prob:bound}), and more explicitly by $K(\mu)\ln 2 + \kappa$ (\cref{prob:bound_explicit}). Since a convergent sum of non-negative terms has terms tending to zero, the per-step error $d_t \to 0$. In words: \emph{on average over histories, $M$'s next-symbol predictions converge to $\mu$'s.}

\textbf{Roadmap.}
\begin{itemize}\itemsep=2pt
  \item \cref{prob:chain_rule} (\emph{Chain rule}): KL divergence of a pair splits into a marginal plus a conditional term.
  \item \cref{prob:telescope} (\emph{Telescoping}): iterating the chain rule turns a joint KL over a whole history into a sum of per-step conditional KLs.
  \item \cref{prob:pinsker} (\emph{Pinsker}): squared error is dominated by KL divergence---this converts the quantity we care about into one the chain rule can handle.
  \item \cref{prob:bound} (\emph{Main bound}): assemble the above and use the mixture domination $M \geq w_\mu\,\mu$ to conclude $S_\infty^\mu \leq -\ln w_\mu$.
  \item \cref{prob:bound_explicit} (\emph{Explicit bound}): refine this to $K(\mu)\ln 2 + \kappa$ using the universal prior.
\end{itemize}

%===================================
\section*{Setup and Definitions}

For more setup and definitions, see \href{https://www.lesswrong.com/posts/HSDumToH57nSRdLST/a-technical-introduction-to-solomonoff-induction-without-k}{the corresponding post on Solomonoff induction}.

\begin{defbox}[Notation and symbols]
\begin{itemize}\itemsep=2pt
  \item $\mathbb{B} = \{\texttt{0}, \texttt{1}\}$ --- the binary alphabet
  \item $\mathbb{B}^*$ --- set of all finite binary strings
  \item $\epsilon$ --- the empty string
  \item $x_{1:n} := x_1 x_2 \cdots x_n$ --- a finite binary string of length $n$
  \item $x_{<t} := x_1 x_2 \cdots x_{t-1}$ --- string prefix up to (but not including) position $t$
  \item $X_{1:n} := X_1, \dots, X_n$ and $X_{<t} := X_1, \dots, X_{t-1}$ --- the corresponding random variables
\end{itemize}
\end{defbox}

\begin{definition}[(Semi)probability distribution]\label{def:semiprob}
A \textbf{semiprobability distribution} over a finite set $\mathcal{X}$ is a function $Q: \mathcal{X} \to [0, 1]$ with $Q(x) \geq 0$ and $\sum_{x \in \mathcal{X}} Q(x) \leq 1$. A semiprobability distribution is a \textbf{probability distribution} if $\sum_{x \in \mathcal{X}} Q(x) = 1$.
\end{definition}

\begin{definition}[(Semi)measure]\label{def:semimeasure}
A \textbf{semimeasure} over infinite binary sequences is a function $\nu: \mathbb{B}^* \to [0, 1]$ with $\nu(\epsilon) \leq 1$ and $\nu(x) \geq \nu(x\texttt{0}) + \nu(x\texttt{1})$ for all $x \in \mathbb{B}^*$.
A semimeasure is a \textbf{measure} if equality holds for both equations.
\end{definition}

\begin{definition}[Model class, prior, and Solomonoff mixture $M$]\label{def:mixture}
Let $\Msol = \{\nu_1, \nu_2, \dots\}$ be the countable class of semimeasures, enumerated as $\nu_1, \nu_2, \dots$, and let $w_\nu > 0$ be a prior weight on each $\nu \in \Msol$ with $\sum_{\nu \in \Msol} w_\nu \leq 1$. We assume the true measure satisfies $\mu \in \Msol$.

The \textbf{Solomonoff mixture} (universal a priori distribution) is
\begin{equation}\label{eq:mixture_expression}
M(x) ~:=~ \sum_{\nu \in \Msol} w_\nu\, \nu(x)
~\geq~ w_\mu\, \mu(x),
\end{equation}
the inequality holding because $\mu \in \Msol$ contributes a single non-negative term to the sum. As a mixture of semimeasures, $M$ is itself a semimeasure.
\end{definition}

\begin{defbox}[The universal prior]
The canonical choice of prior is the \textbf{Solomonoff prior} $w_i := 2^{-K(i)}$, where $K(i)$ is the Kolmogorov complexity of index $i$ (the length of the shortest program that outputs $i$). For a measure $\mu$, write $K(\mu) := \min\{ K(i) : \nu_i = \mu \}$.

The Solomonoff prior is \emph{dominant}: there is a constant $c > 0$ (independent of the sequence) such that
\begin{equation}\label{eq:second:representation}
M(x) ~>~ c \cdot \sum_{i} w_i\, \nu_i(x).
\end{equation}
\end{defbox}

%===================================
\section{Chain rule for KL divergence}\label{prob:chain_rule}

\begin{definition}[Kullback--Leibler divergence]\label{def:kl_divergence}
Let $P$ be a probability distribution and $Q$ a \textbf{semi}probability distribution over a finite set $\mathcal{X}$. The \textbf{Kullback--Leibler (KL) divergence} from $P$ to $Q$ is
\[
D_{\text{KL}}(P \,\|\, Q) := \sum_{x \in \mathcal{X}} P(x) \ln \frac{P(x)}{Q(x)},
\]
where we define $0 \ln \frac{0}{q} := 0$ for any $q \geq 0$, and $p \ln \frac{p}{0} := \infty$ for $p > 0$.
\end{definition}

\begin{definition}[Joint and conditional KL divergence]\label{def:joint_kl_divergence}
Let $P$ be a joint probability distribution and $Q$ a joint \textbf{semi}probability distribution over random variables $X, Y$ taking values in $\mathcal{X}, \mathcal{Y}$. The \textbf{joint} and \textbf{conditional} KL divergences are
\[
D_{\text{KL}}(P(X,Y) \,\|\, Q(X,Y)) := \sum_{x, y} P(x,y) \ln \frac{P(x,y)}{Q(x,y)},
\qquad
D_{\text{KL}}(P(Y \mid X) \,\|\, Q(Y \mid X)) := \sum_{x} P(x) \sum_{y} P(y \mid x) \ln \frac{P(y \mid x)}{Q(y \mid x)},
\]
with the same conventions as \cref{def:kl_divergence}.
\end{definition}

\bigskip

\noindent
\textbf{[10]} We show that KL divergence satisfies the chain rule:
\[
D_{\text{KL}}(P(X,Y) \,\|\, Q(X,Y))
=
D_{\text{KL}}(P(X) \,\|\, Q(X))
+
D_{\text{KL}}(P(Y \mid X) \,\|\, Q(Y \mid X)).
\]

\begin{solution}
We expand the joint KL divergence using the standard chain rule for probabilities, $P(x,y) = P(x) P(y \mid x)$:
\begin{align*}
& D_{\text{KL}}(P(X,Y) \,\|\, Q(X,Y)) \\
&= \sum_{x,y} P(x,y)
\ln \frac{P(x,y)}{Q(x,y)} \\
&= \sum_{x,y} P(x) P(y \mid x)
\ln \frac{P(x) P(y \mid x)}{Q(x) Q(y \mid x)} \\
&= \sum_{x} \sum_y P(x) P(y \mid x)
\left[
\ln \frac{P(x)}{Q(x)}
+
\ln \frac{P(y \mid x)}{Q(y \mid x)}
\right] \\
&= \sum_x P(x) \ln \frac{P(x)}{Q(x)}
\underbrace{\sum_y P(y \mid x)}_{=1}
 + \sum_x P(x) \sum_y P(y \mid x)
\ln \frac{P(y \mid x)}{Q(y \mid x)} \\
&= D_{\text{KL}}(P(X) \,\|\, Q(X))
+
D_{\text{KL}}(P(Y \mid X) \,\|\, Q(Y \mid X)).
\qedhere
\end{align*}
\end{solution}

%===================================
\section{Telescoping the chain rule}\label{prob:telescope}

The two-variable chain rule of \cref{prob:chain_rule} extends, by induction, to a sequence of variables. The point of this telescoping identity is that it lets us trade a single joint KL divergence over an entire history for a sum of per-step conditional KL divergences (and vice versa)---exactly the bridge we will need in \cref{prob:bound}.

\bigskip

\noindent
\textbf{[10]} Let $P$ be a joint probability distribution over a sequence of random variables $X_{1:n} := X_1, \dots, X_n$ taking values in $\mathcal{X}_1 \times \cdots \times \mathcal{X}_n$, and let $Q$ be a \textbf{semi}probability distribution over the same variables. Show by induction on $n$ (using \cref{prob:chain_rule}) that
\[
D_{\text{KL}}(P(X_{1:n}) \,\|\, Q(X_{1:n})) = \sum_{t=1}^{n} D_{\text{KL}}(P(X_t \mid X_{<t}) \,\|\, Q(X_t \mid X_{<t})),
\]
where $P(X_1 \mid X_{<1}) := P(X_1)$ and similarly for $Q$.

\begin{solution}
We prove this by induction on $n$.

\textbf{Base case ($n = 1$):} The sum reduces to $D_{\text{KL}}(P(X_1 \mid X_{<1}) \,\|\, Q(X_1 \mid X_{<1})) = D_{\text{KL}}(P(X_1) \,\|\, Q(X_1))$, which is just the left-hand side.

\textbf{Induction step:} Apply the chain rule from \cref{prob:chain_rule} with $X = X_{<n}$ and $Y = X_n$:
\begin{align*}
D_{\text{KL}}(P(X_{1:n}) \,\|\, Q(X_{1:n}))
&= D_{\text{KL}}(P(X_{<n}) \,\|\, Q(X_{<n})) + D_{\text{KL}}(P(X_n \mid X_{<n}) \,\|\, Q(X_n \mid X_{<n})).
\end{align*}
By the induction hypothesis, the first term satisfies
\[
D_{\text{KL}}(P(X_{<n}) \,\|\, Q(X_{<n})) = \sum_{t=1}^{n-1} D_{\text{KL}}(P(X_t \mid X_{<t}) \,\|\, Q(X_t \mid X_{<t})).
\]
Combining both gives
\[
D_{\text{KL}}(P(X_{1:n}) \,\|\, Q(X_{1:n})) = \sum_{t=1}^{n} D_{\text{KL}}(P(X_t \mid X_{<t}) \,\|\, Q(X_t \mid X_{<t})). \qedhere
\]
\end{solution}

%===================================
\section{Pinsker's inequality}\label{prob:pinsker}

The next ingredient relates squared prediction error to KL divergence. We state the form we will actually use---the general \emph{semiprobability} version---and prove its proof in \cref{app:pinsker}. The exercise here is the key analytic step: the binary inequality.

\begin{theorem}[Pinsker's inequality for semiprobability distributions]\label{thm:pinsker}
Let $P$ be a probability distribution over $\mathbb{B} = \{\texttt{0}, \texttt{1}\}$
and let $Q$ be a semiprobability distribution over $\mathbb{B}$, i.e., $q_0, q_1 \geq 0$
with $q_0 + q_1 \leq 1$. Write $p_x = P(x)$ and $q_x = Q(x)$. Then
\[
    \sum_{x \in \mathbb{B}} \big( q_x - p_x \big)^2 \;\leq\; \sum_{x \in \mathbb{B}} p_x \ln \frac{p_x}{q_x}.
    \tag{$\star$}
\]
The proof is deferred to \cref{app:pinsker}; it reduces the general case to the binary inequality of \cref{lem:pinskers_binary_inequality} below.
\end{theorem}

\begin{lemma}[Pinsker's binary inequality]\label{lem:pinskers_binary_inequality}
For $0 \leq p \leq 1$ and $0 < q < 1$ we have
\[
    2(p-q)^2 ~\leq~ p \ln \frac{p}{q} + (1-p) \ln \frac{1-p}{1-q}.
\]
\end{lemma}

\noindent
\textbf{[15]} Prove \cref{lem:pinskers_binary_inequality}.

\textit{Hint:} Define a function $f(x)$ such that the right-hand side can be written as $f(p) - f(q)$ (try seperating the LHS into terms with and without $q$). Then use the fundamental theorem of calculus, $f(p) - f(q) = \int_q^p f'(x) \, dx$, together with the inequality $x(1-x) \leq 1/4$ for $0 \leq x \leq 1$, to obtain the result.

\begin{solution}
Define $f(x) := p \ln x + (1-p) \ln (1-x)$, so that
\[
p \ln \frac{p}{q} + (1-p) \ln \frac{1-p}{1-q} = f(p) - f(q)
= \int_q^p f'(x)\, dx
= \int_q^p \frac{p-x}{x(1-x)}\, dx,
\]
using $f_p'(x) = \frac{p}{x} - \frac{1-p}{1-x} = \frac{p-x}{x(1-x)}$. Recall also that $x(1-x) \leq \tfrac14$ on $[0,1]$, so $\frac{1}{x(1-x)} \geq 4$.

\textbf{Case $q \leq p$.} For every $x \in [q,p]$ we have $x \leq p$, hence $p - x \geq 0$. Multiplying $\frac{1}{x(1-x)} \geq 4$ by the nonnegative factor $p-x$ preserves the direction:
\[
\frac{p-x}{x(1-x)} \;\geq\; 4(p-x) \qquad \text{for } x \in [q,p].
\]
Integrating over $[q,p]$,
\[
\int_q^p \frac{p-x}{x(1-x)}\, dx \;\geq\; 4 \int_q^p (p-x)\, dx \;=\; 2(p-q)^2.
\]

\textbf{Case $q > p$.} This is the \emph{same} computation with two sign reversals that cancel, so we do not repeat the algebra. First, reversing the limits, $\int_q^p = -\int_p^q$. Second, on $[p,q]$ we now have $x \geq p$, so $p-x \leq 0$, and multiplying $\frac{1}{x(1-x)} \geq 4$ by this nonpositive factor \emph{reverses} the pointwise inequality to $\frac{p-x}{x(1-x)} \leq 4(p-x)$. Integrating over $[p,q]$ and negating both sides (the second flip) restores the original direction:
\[
\int_p^q \frac{p-x}{x(1-x)}\, dx \;\leq\; 4\int_p^q (p-x)\, dx
\;\;\Longrightarrow\;\;
\int_q^p \frac{p-x}{x(1-x)}\, dx \;\geq\; 4\int_q^p (p-x)\, dx \;=\; 2(p-q)^2.
\]
Either way, $f_p(p) - f_p(q) \geq 2(p-q)^2$. \qedhere

\medskip
\emph{Remark (one-shot version).} Equivalently, note $f_p'(x) - 4(p-x) = (p-x)\big[\tfrac{1}{x(1-x)} - 4\big]$ has the same sign as $p-x$, since the bracket is $\geq 0$ on $(0,1)$. Hence the integrand is $\geq 0$ on $[q,p]$ when $q \leq p$ and $\leq 0$ on $[p,q]$ when $q > p$, so $\int_q^p \big[f_p'(x) - 4(p-x)\big]\, dx \geq 0$ in both cases.

\medskip
\emph{Alternative (no case split).} The same FTC computation avoids the case analysis entirely via a linear change of variables along the segment from $q$ to $p$. Using $2(p-q)^2 = 4\int_q^p (p-x)\, dx$,
\[
f_p(p) - f_p(q) - 2(p-q)^2 = \int_q^p \big[\,f_p'(x) - 4(p-x)\,\big]\, dx.
\]
Substitute $x = q + t(p-q)$ for $t \in [0,1]$, so $dx = (p-q)\, dt$ and $p - x = (1-t)(p-q)$:
\[
\int_q^p \big[\,f_p'(x) - 4(p-x)\,\big]\, dx
= \int_0^1 (1-t)\,(p-q)^2 \Big[\tfrac{1}{x(1-x)} - 4\Big]\, dt
\;\geq\; 0,
\]
since $(1-t) \geq 0$, $(p-q)^2 \geq 0$, and $\tfrac{1}{x(1-x)} - 4 \geq 0$ because $x = q + t(p-q) \in [0,1]$ is a convex combination of $q, p \in [0,1]$, so $x(1-x) \leq \tfrac14$. The orientation issue disappears because the $(p-x)$ in the integrand and the $(p-q)$ from $dx$ combine into $(p-q)^2 \geq 0$.
\end{solution}

%===================================
\section{Cumulative prediction error bound}\label{prob:bound}

Let $\mu$ be any measure on binary sequences, and assume it generates our actual universe. Let $M$ be the Solomonoff mixture distribution of \cref{def:mixture}.

Define the cumulative expected prediction error of predicting sequences sampled by $\mu$ via $M$ as
\[
S^{\mu}_{\infty} := \sum_{t = 1}^{\infty} \sum_{x_{<t} \in \mathbb{B}^*} \mu(x_{<t}) \sum_{x_t \in \mathbb{B}} \big( M(x_t \mid x_{<t}) - \mu(x_t \mid x_{<t}) \big)^2.
\]

\noindent
\textbf{[20]} Show that
\[
S_{\infty}^{\mu} \leq - \ln w_\mu.
\]

\textit{Hint:} Apply \cref{thm:pinsker} to each summand to pass from squared error to KL divergence, then make the infinite sum explicit as a limit of finite sums up to $n$, then use \cref{prob:telescope} to telescope the KL divergences, and finally use the mixture representation of $M$.

\begin{solution}
For each $t$ and $x_{<t}$, the conditional $\mu(\cdot \mid x_{<t})$ is a probability distribution over $\mathbb{B}$ (since $\mu$ is a measure), while $M(\cdot \mid x_{<t})$ is a semiprobability distribution over $\mathbb{B}$ (since $M$ is a semimeasure). Thus, by \cref{thm:pinsker}:
\[
\sum_{x_t \in \mathbb{B}} \big( M(x_t \mid x_{<t}) - \mu(x_t \mid x_{<t}) \big)^2 \leq \sum_{x_t \in \mathbb{B}} \mu(x_t \mid x_{<t}) \ln \frac{\mu(x_t \mid x_{<t})}{M(x_t \mid x_{<t})}.
\]
Thus:
\begin{align*}
S_{\infty}^{\mu} &\leq \sum_{t=1}^{\infty} \sum_{x_{<t} \in \mathbb{B}^*} \mu(x_{<t}) \sum_{x_t \in \mathbb{B}} \mu(x_t \mid x_{<t}) \ln \frac{\mu(x_t \mid x_{<t})}{M(x_t \mid x_{<t})} \\
&= \lim_{n \to \infty} \sum_{t=1}^{n} D_{\text{KL}}\big( \mu(X_t \mid X_{<t}) \,\|\, M(X_t \mid X_{<t}) \big),
\end{align*}
where the inner sum over $x_{<t}$ weighted by $\mu(x_{<t})$ gives exactly the conditional KL divergence. By \cref{prob:telescope} (the telescoped chain rule), this equals
\[
S_{\infty}^{\mu} \leq \lim_{n \to \infty} D_{\text{KL}}\big( \mu(X_{1:n}) \,\|\, M(X_{1:n}) \big) = \lim_{n \to \infty} \sum_{x_{1:n} \in \mathbb{B}^n} \mu(x_{1:n}) \ln \frac{\mu(x_{1:n})}{M(x_{1:n})}.
\]
Now, by the mixture representation \eqref{eq:mixture_expression},
\[
M(x_{1:n}) = \sum_{\nu \in \Msol} w_\nu\, \nu(x_{1:n}) \geq w_\mu\, \mu(x_{1:n}).
\]
Therefore $\frac{\mu(x_{1:n})}{M(x_{1:n})} \leq \frac{1}{w_\mu}$, and so
\[
S_{\infty}^{\mu} \leq \lim_{n \to \infty} \sum_{x_{1:n}} \mu(x_{1:n}) \ln \frac{1}{w_\mu} = -\ln w_\mu \cdot \lim_{n \to \infty} \underbrace{\sum_{x_{1:n}} \mu(x_{1:n})}_{= 1} = -\ln w_\mu,
\]
where $\sum_{x_{1:n}} \mu(x_{1:n}) = 1$ since $\mu$ is a measure. \qedhere
\end{solution}

%===================================
\section{Explicit complexity bound}\label{prob:bound_explicit}

\noindent
\textbf{[10]} Under the same setup as \cref{prob:bound}, show that there exists a constant $\kappa \in \mathbb{R}$ (independent of $\mu$) such that
\[
S_{\infty}^{\mu} \leq K(\mu) \ln 2 + \kappa,
\]
where $K(\mu) := \min\{ K(i) : \nu_i = \mu \}$ is the Kolmogorov complexity of $\mu$ (minimized over all indices $i$ in the effective enumeration that give rise to $\mu$).

\textit{Hint:} Follow the same strategy as \cref{prob:bound}, but instead of the mixture representation \eqref{eq:mixture_expression}, use the lower bound \eqref{eq:second:representation} together with $w_i = 2^{-K(i)}$.

\begin{solution}
The proof follows the same steps as \cref{prob:bound} up to the KL divergence expression:
\[
S_{\infty}^{\mu} \leq \lim_{n \to \infty} \sum_{x_{1:n}} \mu(x_{1:n}) \ln \frac{\mu(x_{1:n})}{M(x_{1:n})}.
\]
Now, by \eqref{eq:second:representation}, we have
\[
M(x_{1:n}) > c \cdot \sum_{i} w_i\, \nu_i(x_{1:n}) \geq c \cdot w_{i^*}\, \mu(x_{1:n}),
\]
where $i^*$ is any index with $\nu_{i^*} = \mu$. Therefore $\frac{\mu(x_{1:n})}{M(x_{1:n})} < \frac{1}{c \cdot w_{i^*}}$, and so
\[
S_{\infty}^{\mu} \leq \lim_{n \to \infty} \sum_{x_{1:n}} \mu(x_{1:n}) \ln \frac{1}{c \cdot w_{i^*}} = -\ln c - \ln w_{i^*} = -\ln c + K(i^*) \ln 2,
\]
using $w_{i^*} = 2^{-K(i^*)}$ and $\sum_{x_{1:n}} \mu(x_{1:n}) = 1$. Since this holds for every $i^*$ with $\nu_{i^*} = \mu$, we may take the minimum over all such indices to obtain
\[
S_{\infty}^{\mu} \leq K(\mu) \ln 2 + \kappa,
\]
where $\kappa = -\ln c$. \qedhere
\end{solution}

\leon{Throughout, there are some subtle trivial compatibilities between semimeasures and semiprobability distributions that I use, e.g. that semimeasures are semiprobability distributions on fixed-length sequences, and that they give rise to CONDITIONALS that are actual probability distributions. Maybe I could mention that or make these their own exercises? Hhm.}

%===================================
\section*{Interpretation}

\cref{prob:bound,prob:bound_explicit} show that the \emph{cumulative} expected prediction error $S_\infty^\mu$ is finite. Since $S_\infty^\mu$ is a sum of non-negative terms, this immediately implies that the per-step expected prediction error
\[
d_t := \sum_{x_{<t} \in \mathbb{B}^*} \mu(x_{<t}) \sum_{x_t \in \mathbb{B}} \big( M(x_t \mid x_{<t}) - \mu(x_t \mid x_{<t}) \big)^2
\]
converges to zero as $t \to \infty$. In other words, $M$'s predictions become indistinguishable from $\mu$'s predictions on average over histories.

%===================================
\clearpage
\appendix
\renewcommand{\thesection}{\Alph{section}}
\crefname{section}{Appendix}{Appendices}
\Crefname{section}{Appendix}{Appendices}

\section{Proof of Pinsker for semiprobability distributions}\label{app:pinsker}

We collect two auxiliary facts, then prove \cref{thm:pinsker} in full generality.

\begin{lemma}[Logarithm bound]\label{lem:log_bound}
For all $z > 0$, \quad $\ln z \geq 1 - \dfrac{1}{z}$.
\end{lemma}

\begin{proof}
Define $g(z) = \ln z - 1 + \frac{1}{z}$. We show $g(z) \geq 0$ for all $z > 0$. We have
\[
g'(z) = \frac{1}{z} - \frac{1}{z^2} = \frac{z - 1}{z^2},
\]
so $g'(z) < 0$ for $z \in (0, 1)$ and $g'(z) > 0$ for $z > 1$; hence $g$ has a global minimum at $z = 1$. Since $g(1) = \ln 1 - 1 + 1 = 0$, we conclude $g(z) \geq 0$, i.e.\ $\ln z \geq 1 - \frac{1}{z}$.
\end{proof}

\begin{fact}[AM--GM inequality]\label{fact:amgm}
For all $a, b \geq 0$, \quad $a + b \geq 2\sqrt{ab}$.
\end{fact}

\begin{proof}[Proof of \cref{thm:pinsker}]
We handle the boundary case $q_x = 0$ first, then reduce the interior case to
\cref{lem:pinskers_binary_inequality}.

\medskip
\textbf{Case 1: $q_x = 0$ for some $x \in \mathbb{B}$.}
If $p_x > 0$, then the right-hand side of $(\star)$ contains
$p_x \ln \tfrac{p_x}{0} = +\infty$, so $(\star)$ holds trivially.
If instead $p_x = 0$, that atom contributes $0$ to both sides
(using the convention $0 \ln \tfrac{0}{0} = 0$). The other atom $x'$ then
satisfies $p_{x'} = 1$ and $q_{x'} \in [0, 1]$, so $(\star)$ reduces to
\[
    (1 - q_{x'})^2 \;\leq\; \ln \tfrac{1}{q_{x'}}.
\]
By \cref{lem:log_bound} applied to $z = 1/q_{x'}$, we have
$\ln \tfrac{1}{q_{x'}} \geq 1 - q_{x'}$, and since $1 - q_{x'} \in [0, 1]$
we get $(1 - q_{x'})^2 \leq 1 - q_{x'} \leq \ln \tfrac{1}{q_{x'}}$.

\medskip
\textbf{Case 2: $q_0, q_1 > 0$.}
Define
\[
    \Phi(q_0, q_1) \;:=\; \sum_{x} p_x \ln \tfrac{p_x}{q_x} - \sum_{x} (q_x - p_x)^2,
\]
so that $(\star)$ is equivalent to $\Phi(q_0, q_1) \geq 0$.
We show $\Phi$ is decreasing in $q_1$ (for fixed $q_0$):
\[
    \frac{\partial \Phi}{\partial q_1}
    \;=\; -\frac{p_1}{q_1} - 2(q_1 - p_1)
    \;=\; 2p_1 - \Big(\tfrac{p_1}{q_1} + 2q_1\Big).
\]
By \cref{fact:amgm} applied to $a = p_1/q_1$ and $b = 2q_1$,
\[
    \tfrac{p_1}{q_1} + 2q_1 \;\geq\; 2\sqrt{2p_1},
\]
hence
\[
    \frac{\partial \Phi}{\partial q_1} \;\leq\; 2p_1 - 2\sqrt{2p_1}
    \;=\; 2\sqrt{p_1}\big(\sqrt{p_1} - \sqrt{2}\big) \;<\; 0,
\]
using $p_1 \leq 1 < 2$. Therefore, extending $q_1$ from its current value
up to $1 - q_0$ only decreases $\Phi$:
\[
    \Phi(q_0, q_1) \;\geq\; \Phi(q_0, 1 - q_0).
\]
But $\Phi(q_0, 1 - q_0) \geq 0$ is exactly \cref{lem:pinskers_binary_inequality}
applied with $p = p_0$ and $q = q_0$.
\end{proof}


\section{Knuth's Difficulty Scale}\label{sec:appendix_difficulty}

Each problem carries a difficulty rating in square brackets, following Knuth's rating scheme for exercises in slightly adapted form. The rating assumes that the material in the preceding problems (on which the problem depends) has been understood. In-between values are possible.

\begin{itemize}\itemsep=2pt
    \item[\textbf{[00]}] \emph{Very easy.} Solvable from the top of your head.
    \item[\textbf{[10]}] \emph{Easy.} Needs 15 minutes to think, possibly pencil and paper.
    \item[\textbf{[20]}] \emph{Average.} May take 1--2 hours to answer completely.
    \item[\textbf{[30]}] \emph{Moderately difficult or lengthy.} May take several hours to a day.
    \item[\textbf{[40]}] \emph{Quite difficult or lengthy.} Often a significant research result.
    \item[\textbf{[50]}] \emph{Open research problem.} An obtained solution should be published.
\end{itemize}

\end{document}
